% ============================================================================
%  Depuración en hardware
% ============================================================================
\section{Depuración en hardware}
\label{sec:depuracion}

La puesta en marcha en la tarjeta presentó fallas que, a primera vista,
parecían errores de software. Cada una se aisló cambiando una sola variable a
la vez: otra tarjeta, otro código, otra fuente o una medición en un punto
concreto. La tabla~\ref{tab:depuracion} resume el proceso en el orden en que
ocurrió.

\begin{table}[H]
  \centering
  \caption{Fallas encontradas, diagnóstico y solución.}
  \label{tab:depuracion}
  \footnotesize
  \begin{tabularx}{\textwidth}{@{}c >{\raggedright\arraybackslash}p{2.9cm} L L@{}}
    \toprule
    \textbf{\#} & \textbf{Síntoma} & \textbf{Diagnóstico} & \textbf{Causa y solución} \\
    \midrule
    1 & El código de la guía no compila. &
      Al copiarlo del PDF quedaron los encabezados de página dentro del código. &
      Se retiraron los encabezados y se movió la declaración de \texttt{velocidad} al inicio de \texttt{main()} (compatibilidad C90). \\
    2 & La aplicación del PICkit se bloquea al elegir la familia del PIC. &
      Ocurre durante la detección automática del chip con un programador clon. &
      \emph{Manual Device Select} y selección manual del PIC16F877A. \\
    3 & \emph{Programming failed at address 0x000000} en una tarjeta. &
      La misma grabación funciona en otra tarjeta: el problema es de esa tarjeta. Antes se había conectado el ICSP invertido. &
      VPP (\SI{13}{\volt}) aplicado en RB3/PGM. PIC dañado; se trabaja en la otra tarjeta. \\
    4 & Los motores no se mueven con el código del docente. &
      EEP medido en \SI{0.14}{\volt}; el esquemático muestra EEP sin conexión. &
      Driver dormido. Se habilita EEP por hardware (sección~\ref{sec:eep}). \\
    5 & Con el código propio los motores se mueven a ratos, con demoras, o solo uno. &
      EEP flotante: el ruido del PWM lo despertaba por instantes. A baja velocidad, además, la zona muerta del motor. &
      Habilitación fija de EEP y compensación de zona muerta ($V_{MIN} = 300$). \\
    6 & En Dabble, $\triangle$ hace retroceder el carro; ninguna dirección responde bien. &
      Con los LEDs de diagnóstico: solo cambia LED2, nunca LED1. Los bytes llegan alterados. &
      HC-05 a una velocidad distinta de 9600 baudios. Detección automática de baudios (sección~\ref{sec:autobaud}). \\
    7 & Tras un comando, el carro queda girando y no obedece. &
      Dabble muestra \emph{HC-05 disconnected} justo al arrancar el motor. &
      La pila de \SI{9}{\volt} cae y el HC-05 se reinicia; el PIC conservaba la última orden. Freno por pérdida de señal (sección~\ref{sec:failsafe}). \\
    8 & Tras agregar el freno de seguridad, el motor no arranca. &
      El freno se activaba con un único error de trama causado por el ruido del motor. &
      Se exigen tres errores consecutivos para frenar. \\
    9 & Los tres LEDs parpadean al presionar una dirección. &
      Es la señal de arranque del PIC: hubo un reinicio por bajo voltaje (BOR). Sin motores no ocurre. &
      Fuente insuficiente. Se requieren baterías 18650 o una fuente de laboratorio. \\
    10 & Con EEP conectado a \SI{5}{\volt} mediante \SI{10}{\kilo\ohm}, el motor responde a ratos. &
      EEP mide \SI{1.41}{\volt}, en la zona no definida de nSLEEP, y a veces \SI{5}{\volt}. &
      Divisor con el \emph{pull-down} del módulo y contacto intermitente. Cable directo a \SI{5}{\volt}, sin continuidad con OUT1. \\
    \bottomrule
  \end{tabularx}
\end{table}

\subsection{Lecciones}

\begin{itemize}
  \item \textbf{Los LEDs indican órdenes, no resultados.} Que LED3 se encienda
        demuestra que el PIC ordena mover el motor, no que el motor gire. Para
        separar ambas cosas hubo que medir en la cadena PIC $\rightarrow$
        driver $\rightarrow$ motor.
  \item \textbf{Una entrada de habilitación nunca debe quedar flotando.} Un
        nivel indefinido produce fallas intermitentes, que son las más
        difíciles de diagnosticar.
  \item \textbf{La alimentación se prueba primero.} Varias fallas que parecían
        de comunicación o de código eran caídas de voltaje al arrancar los
        motores.
  \item \textbf{Un mecanismo de protección también puede fallar.} El freno por
        pérdida de señal, demasiado sensible, detenía el carro por ruido; hubo
        que ajustarlo a las condiciones reales.
  \item \textbf{Verificar la orientación del programador antes de cualquier
        clic.} La aplicación aplica voltaje apenas detecta el chip.
\end{itemize}
